\chapter{Cadre conceptuel et méthodologique}

\section{Introduction}
La PAGe Platform repose sur un cadre conceptuel et des règles de calcul qui ont été définis par notre encadrant, Pr. B. BOUNABAT. Ce cadre établit les concepts (risque, indicateur, mécanisme, pilier, maturité) et les formules d'agrégation utilisées dans les trois modules de la plateforme.

Notre rôle n'a pas été de concevoir ces formules, mais de les comprendre et de les implémenter dans l'application. Ce chapitre les présente telles qu'elles nous ont été fournies, afin que le lecteur puisse comprendre la logique de calcul derrière chaque score affiché par la plateforme.

% ================================================================
\section{Concepts fondamentaux}

\subsection{Risque algorithmique}
Un risque algorithmique désigne tout effet indésirable pouvant résulter de l'utilisation d'un algorithme dans un processus décisionnel d'une administration publique. Exemples : biais discriminatoires, opacité des décisions, fragmentation de l'information entre administrations.

Dans la plateforme, chaque risque est une unité analytique autonome, identifiée et suivie indépendamment.

\subsection{Indicateur}
Un indicateur est une variable quantifiable qui mesure un aspect d'un risque. Par exemple : le taux d'intégration des bases de données, le nombre de recours contre des décisions automatisées, le niveau de documentation des modèles.

Chaque indicateur possède :
\begin{itemize}
    \item une \textbf{valeur brute} ($Val_{Ind_i}$) : renseignée par l'utilisateur ;
    \item un \textbf{poids} ($W_{Ind_i}$) : importance relative, avec $\sum W_{Ind_i} = 1$ ;
    \item un \textbf{statut} : POSITIVE, NEGATIVE ;
    \item une \textbf{plage} : valeur min et valeur max.
\end{itemize}

\subsection{Mécanisme de mitigation (RMM)}
Un RMM (\textit{Risk Mitigation Mechanism}) est un levier concret pour réduire un risque : audit algorithmique indépendant, interopérabilité technique, encadrement juridique, formation des agents, etc.

Chaque mécanisme est évalué selon sa capacité réelle d'activation, qui dépend de la préparation de l'organisation sur différents piliers.

\subsection{Piliers de gouvernance (Key Pillars)}
Les piliers représentent les fondements institutionnels nécessaires à l'activation des mécanismes. La plateforme en définit six :
\begin{itemize}
    \item \textbf{Governance} : cadre décisionnel et stratégique ;
    \item \textbf{Organizational} : capacités organisationnelles ;
    \item \textbf{Legal} : cadre juridique et réglementaire ;
    \item \textbf{Technical} : infrastructure technique ;
    \item \textbf{Human} : compétences humaines ;
    \item \textbf{Financial} : ressources financières.
\end{itemize}

Chaque pilier est évalué via un \textit{Readiness Level} (RL), avec $RL_i \in [0, 1]$.

\subsection{Dimensions, Facteurs et Items}
Chaque pilier est décomposé en :
\begin{itemize}
    \item \textbf{Dimensions} : axes opérationnels ;
    \item \textbf{Facteurs} : sous-décomposition des dimensions ;
    \item \textbf{Items} : questions concrètes posées via un questionnaire, avec des réponses de 1 à 5.
\end{itemize}

% ================================================================
\section{Module O-PAGe : mesure de la gravité des risques}

\subsection{Normalisation des indicateurs}
La valeur normalisée dépend du statut de l'indicateur :

$$
Val^{norm}_{Ind_i} = \begin{cases}
Val_{Ind_i} & \text{si } Statut_{Ind_i} = \text{POSITIVE} \\
1 - Val_{Ind_i} & \text{si } Statut_{Ind_i} = \text{NEGATIVE} \\
\end{cases}
$$

Avec $Val^{norm}_{Ind_i} \in [0, 1]$.

Un indicateur POSITIVE : plus la valeur est élevée, plus le risque est élevé. Un indicateur NEGATIVE : c'est l'inverse.

\subsection{Calcul du Risk Score}
Le score de gravité est l'agrégation pondérée des indicateurs normalisés :

$$
Risk\_Score = \sum_{i} W_{Ind_i} \cdot Val^{norm}_{Ind_i}
$$

Avec $\sum_{i} W_{Ind_i} = 1$. Le résultat est entre 0 et 1.

\subsection{Catégorisation}
Le score est traduit en niveaux :

\begin{table}[h!]
\centering
\renewcommand{\arraystretch}{1.3}
\begin{tabular}{@{}p{0.20\textwidth}p{0.35\textwidth}p{0.35\textwidth}@{}}
\hline
\textbf{Catégorie} & \textbf{Condition} & \textbf{Signification} \\
\hline
Low & $Risk\_Score < 0.25$ & Exposition limitée \\
\hline
Moderate & $0.25 \leq Risk\_Score < 0.50$ & Vigilance requise \\
\hline
High & $0.50 \leq Risk\_Score < 0.75$ & Action nécessaire \\
\hline
Critical & $Risk\_Score \geq 0.75$ & Situation préoccupante \\
\hline
\end{tabular}
\caption{Catégorisation du Risk Score}
\end{table}

\subsection{Exemple}
Un risque avec 3 indicateurs :

\begin{table}[h!]
\centering
\renewcommand{\arraystretch}{1.3}
\begin{tabular}{@{}lcccc@{}}
\hline
\textbf{Indicateur} & \textbf{Valeur brute} & \textbf{Statut} & \textbf{Val. norm.} & \textbf{Poids} \\
\hline
Taux d'erreur & 0.60 & POSITIVE & 0.60 & 0.40 \\
Taux de conformité & 0.80 & NEGATIVE & 0.20 & 0.35 \\
Opacité & 0.50 & POSITIVE & 0.50 & 0.25 \\
\hline
\end{tabular}
\caption{Exemple d'indicateurs}
\end{table}

$$
Risk\_Score = (0.40 \times 0.60) + (0.35 \times 0.20) + (0.25 \times 0.50) = 0.435
$$

Résultat : \textbf{Moderate}.

% ================================================================
\section{Module M-PAGe : mesure de la capacité de mitigation}

\subsection{Agrégation multiniveau}
Les réponses aux items sont de 1 à 5 : $Val_{Item} \in \{1, 2, 3, 4, 5\}$

L'agrégation se fait en plusieurs étapes :

\vspace{0.3cm}

oindent\textbf{Score facteur normalisé :}
$$
Val\_Fact\_Norm = \frac{Avg(Val_{Item}) - 1}{4}
$$

oindent\textbf{Score dimension :}
$$
Val\_Dim = Avg(Val\_Fact\_Norm)
$$

oindent\textbf{Readiness Level du pilier :}
$$
RL_i = Avg(Val\_Dim)
$$

\subsection{Capacité d'un mécanisme (RMMC)}
$$
RMMC_j = \sum_{i} W_{KP_{ji}} \cdot RL_i
$$

Avec $\sum_{i} W_{KP_{ji}} = 1$.

\subsection{Capacité par risque (RMC)}
$$
RMC_k = \sum_{j} W_{RMM_{kj}} \cdot RMMC_j
$$

\subsection{Maturité globale (GPM)}
$$
GPM = \sum_{k} W_{Risk_k} \cdot RMC_k
$$

Le GPM synthétise la capacité globale de l'organisation face à l'ensemble de ses risques.

\subsection{Exemple}
Un mécanisme qui dépend de 3 piliers :

\begin{table}[h!]
\centering
\renewcommand{\arraystretch}{1.3}
\begin{tabular}{@{}lccc@{}}
\hline
\textbf{Pilier} & \textbf{RL} & \textbf{Poids} & \textbf{Contribution} \\
\hline
Governance & 0.70 & 0.40 & 0.280 \\
Technical & 0.55 & 0.35 & 0.193 \\
Human & 0.60 & 0.25 & 0.150 \\
\hline
\multicolumn{3}{@{}l}{\textbf{RMMC}} & \textbf{0.623} \\
\hline
\end{tabular}
\caption{Exemple de calcul RMMC}
\end{table}

% ================================================================
\section{Module I-PAGe : simulation d'impact}

\subsection{Matrice d'impact}
La relation entre mécanismes et indicateurs est formalisée par :
$$
Mat\_RMM\_Ind = M\_Impact
$$

Les lignes = RMM, les colonnes = indicateurs. Chaque valeur $M\_Impact_{ijk}$ traduit l'influence du mécanisme $j$ sur l'indicateur $k$ du risque $i$.

\subsection{Efficacité réelle}
$$
Effi\_RMM_{ij} = RMMC_{ij} \times RMMD_{ij}
$$

C'est le produit de la capacité intrinsèque par le niveau de déploiement réel. Un mécanisme non déployé n'a aucun effet.

\subsection{Évolution des indicateurs}
$$
Ind_{ik}(t+1) = Ind_{ik}(t) - \sum_{j} M\_Impact_{ijk} \cdot Effi\_RMM_{ij}
$$

Chaque mécanisme déployé réduit certains indicateurs. On recalcule ensuite le $Risk\_Score(t+1)$ pour comparer avant/après.

\subsection{Exemple}
Un indicateur à 0.60, un mécanisme avec $M\_Impact = 0.3$ et $Effi\_RMM = 0.5$ :
$$
Ind(t+1) = 0.60 - (0.3 \times 0.5) = 0.45
$$

% ================================================================
\section{Synthèse}

\begin{table}[h!]
\centering
\renewcommand{\arraystretch}{1.4}
\begin{tabular}{@{}p{0.12\textwidth}p{0.25\textwidth}p{0.53\textwidth}@{}}
\hline
\textbf{Module} & \textbf{Score} & \textbf{Formule} \\
\hline
O-PAGe & Risk Score & $\sum W_{Ind_i} \cdot Val^{norm}_{Ind_i}$ \\
\hline
M-PAGe & RL & $Avg(Val\_Dim)$ \\
 & RMMC & $\sum W_{KP_{ji}} \cdot RL_i$ \\
 & RMC & $\sum W_{RMM_{kj}} \cdot RMMC_j$ \\
 & GPM & $\sum W_{Risk_k} \cdot RMC_k$ \\
\hline
I-PAGe & $Ind(t+1)$ & $Ind(t) - \sum M\_Impact \cdot Effi\_RMM$ \\
\hline
\end{tabular}
\caption{Récapitulatif des formules}
\end{table}

\section{Conclusion}
Ce chapitre a présenté les concepts et formules de calcul qui fondent la PAGe Platform. Ces règles, définies par l'encadrant, ont été directement implémentées dans le code de l'application. Le chapitre suivant traite de l'analyse des besoins et de la conception.
